\section{Introduction}
\label{sec:introduction}

During earthquakes, soil elements are subjected to three-dimensional cyclic loading in which the shear stress direction changes continuously as seismic waves propagate through the ground. In practice, however, liquefaction assessment is predominantly based on unidirectional cyclic loading tests, raising concerns that liquefaction resistance may be overestimated when multidirectional effects are neglected.

At the physical-model scale, \citet{Pyke1975} found that multidirectional shaking of dry sand on a shaking table induces greater settlement and requires less stress to cause liquefaction than unidirectional shaking. Centrifuge experiments by \citet{SuLi2008} confirmed that multidirectional shaking accelerates excess pore water pressure build-up, and \citet{ElShafee2017} reported that a 40\% increase in uniaxial shaking amplitude was required to produce a similar pore pressure response to that of biaxial shaking---far exceeding the conventional 10\% correction factor. These model-scale studies established that multidirectional loading is more severe than unidirectional loading, but the coupling of wave propagation, boundary effects, and soil heterogeneity in physical models makes it difficult to isolate the stress-path-level mechanisms responsible for the observed differences.

Element-level tests have provided more direct insight. \citet{IshiharaYamazaki1980} showed in multidirectional simple shear tests that equal-amplitude orthogonal loading requires only approximately 70\% of the cyclic stress ratio needed under unidirectional shear to produce equivalent shear strains. \citet{Kammerer2005} conducted simple shear tests with linear, oval, circular, and figure-8 loading paths, demonstrating that rotation of the shear stress direction promotes faster pore pressure accumulation. \citet{JinGuo2021} further parameterized multidirectional loading by the phase difference $\theta$ between orthogonal shear components in simple shear tests on Toyoura sand, reporting a non-monotonic dependence of liquefaction resistance on $\theta$ and an overall reduction to approximately 70\% of the unidirectional value.

The discrete element method (DEM), originally proposed by \citet{Cundall1979}, has become a widely used tool for investigating liquefaction at the particle scale. \citet{Kuhn2014} demonstrated that three-dimensional DEM assemblies can reproduce the essential features of undrained cyclic liquefaction, including progressive pore pressure build-up and cyclic mobility. \citet{Martin2020} identified micromechanical criteria for the onset of liquefaction through cyclic three-dimensional DEM simulations. \citet{Rahman2021} extended the critical state soil mechanics framework to cyclic liquefaction and post-liquefaction behavior through DEM simulations of triaxial cyclic loading.

Building on this foundation, several DEM studies have compared the liquefaction response under different multidirectional stress paths. \citet{Wei2020}, \citet{WuYang2021servo}, and \citet{Yang2022} systematically varied the path shape, including straight-line, oval, circular, figure-8, and teardrop configurations, and consistently found that path shape significantly influences the number of cycles to liquefaction, with fabric anisotropy evolving in a path-dependent manner. \citet{WuYang2021servo} additionally proposed an undrained servomechanism for imposing arbitrary three-dimensional stress paths and introduced a polyhedral specimen to suppress boundary arching effects.

Other studies have explored additional factors within the multidirectional framework. \citet{JiangXu2022} introduced initial static shear and found that it reduces liquefaction resistance under bidirectional shearing, even though it may enhance resistance under unidirectional loading. \citet{ZhangYang2021} combined bidirectional simple shear experiments with DEM simulations and confirmed good agreement between the two, with coordination number and fabric anisotropy governing the path-dependent response. \citet{XuPan2025} demonstrated that particle shape irregularity increases liquefaction resistance under multidirectional loading, with the effect varying across stress path shapes. \citet{YangHuang2023} and \citet{YangHuangJin2025} extended the DEM approach to reliquefaction, investigating how inherent fabric and multidirectional loading history influence reliquefaction behavior at the microscale.

While these studies have substantially advanced the understanding of path-dependent liquefaction behavior, isolating the specific role of shear stress directionality remains challenging because conventionally defined loading paths differ simultaneously in several respects beyond direction alone:
\begin{enumerate}
  \item Unequal instantaneous shear stress magnitude: The maximum cyclic shear stress ratio (CSR) is typically kept constant across loading paths, but the instantaneous resultant shear stress differs throughout each cycle. Under this convention, straight-line bidirectional paths \citep{WuYang2021servo,JiangXu2022,XuPan2025} yield a peak resultant stress $\sqrt{2}$ times that of the unidirectional path, while circular paths \citep{Wei2020,Yang2022} maintain a constant non-zero resultant equal to the unidirectional peak, never passing through zero.
  \item Unequal stress variation rate: Circular loading imposes a constant rate of stress change, whereas unidirectional loading follows a sinusoidal rate; figure-8 and teardrop paths \citep{WuYang2021servo,XuPan2025} introduce still different temporal patterns. These rate differences are superimposed on the directional change and cannot be separated from it.
\end{enumerate}
These confounding factors make it difficult to isolate the specific influence of shear stress directionality on the liquefaction process.

The present study proposes a double-8 loading method designed to systematically eliminate these confounding factors. The approach is built on three pairwise comparisons. First, an equivalent-magnitude unidirectional path is constructed such that its instantaneous resultant stress magnitude matches that of the single-8 (figure-8) path at every instant, with shear direction being the only difference between the two. Second, a double-8 path is introduced in which the $x$- and $y$-components of the single-8 stress are swapped in alternating cycles, preserving both the stress magnitude and the rate of stress change while introducing periodic axis rotation. By comparing single-8 versus unidirectional, which share the same instantaneous resultant magnitude, the effect of stress magnitude is excluded and the influence of directional change can be assessed. By further comparing double-8 versus single-8, which additionally share the same rate of stress variation, the isolated effect of periodic axis rotation is revealed.

DEM simulations are performed under constant-volume (undrained) conditions using periodic boundaries and ribbed shear walls for shear transmission. The macroscopic liquefaction resistance, cumulative shear work, and strain path development are examined for the three loading types across multiple cyclic stress ratios. Microscopic analyses (mechanical coordination number $\Zm$, fabric anisotropy invariant $\Fc$, and contact density distribution) are then used to explain the particle-scale mechanisms underlying the observed differences in liquefaction resistance.

The remainder of this paper is organized as follows. \Cref{sec:methodology} describes the DEM model, loading path design, calibration, and microscopic descriptors. \Cref{sec:results} presents the macroscopic and microscopic results and discusses the mechanisms governing directionality-dependent liquefaction resistance. \Cref{sec:conclusions} summarizes the main findings.
